% Pages 51–65

% ---- Page 51: Completion of entropy of mixing (problem 2.37 / eq. 2.38) ----

The total entropy of mixing is the sum of the two contributions:
\[
  \Delta S_{\text{mix}} = \Delta S_A + \Delta S_B.
\]
For gases $A$ and $B$ with mole fractions $(1-x)$ and $x$ respectively, expanding each species
into the full volume gives
\begin{align*}
  \Delta S_A &= (1-x)Nk\ln\!\frac{1}{1-x} = -(1-x)Nk\ln(1-x), \\
  \Delta S_B &= xNk\ln\!\frac{1}{x} = -xNk\ln x,
\end{align*}
and therefore
\begin{formula}
\[
  \Delta S_{\text{mix}} = -Nk\bigl[(1-x)\ln(1-x) + x\ln x\bigr].
\]
\end{formula}

It is instructive to derive this result from the combinatorial formula directly. The
$\binom{N}{N_A}$ factor in the multiplicity is essentially the number of distinguishable
configurations of the combined system: it counts the ways to arrange both species that would
have been indistinguishable had all particles been of the same type. Writing $N_A = (1-x)N$ and
$N_B = xN$:
\begin{align*}
  \Delta S_{\text{mixing}}
  &= k\ln\binom{N}{N_A N_B'} = k\ln\!\frac{N!}{(Nx)!\,(N(1-x))!} \\
  &= k\!\left[\ln N! - \ln\bigl((Nx)!\bigr) - \ln\bigl((N(1-x))!\bigr)\right] \\
  &= k\!\left[N\ln N - \cancel{N} - (Nx)\ln(Nx) + \cancel{Nx}
     - N(1-x)\ln\bigl(N(1-x)\bigr) + \cancel{N(1-x)}\right] \\
  &= Nk\!\left[\ln N - x\ln(Nx) - (1-x)\ln\bigl(N(1-x)\bigr)\right] \\
  &= Nk\!\left[\ln N - x(\ln N + \ln x) - (1-x)(\ln N + \ln(1-x))\right] \\
  &= Nk\!\left[-x\ln x - (1-x)\ln(1-x)\right]. \quad\checkmark
\end{align*}

% ---- Page 52: Section 3.1 — Temperature ----

\section{Temperature}

\subsection{How Does Entropy Relate to Other Variables?}

We have defined entropy $S = k\ln\Omega$ and seen that an isolated system tends toward its
highest-entropy macrostate. But entropy is not the most convenient quantity for everyday
measurements. We now ask: how does entropy relate to other macroscopic variables, and in
particular, what is temperature?

Recall that we defined two objects to be at the same \emph{temperature} when they are in
thermal equilibrium with each other. We also know that two objects in thermal equilibrium
attain the maximum total entropy. These two facts together determine the statistical definition
of temperature.

\textbf{Example.} Consider a system of harmonic oscillators: solid $A$ has $N_A = 300$
oscillators, solid $B$ has $N_B = 200$ oscillators, and the total energy is $q_{\text{total}} =
100$ units. The total multiplicity is
\[
  \Omega_{\text{total}} = \binom{300 + q_A}{q_A}\binom{100 - q_A + 200}{100 - q_A}.
\]

At thermal equilibrium we know that
\[
  \pdv{S_{\text{total}}}{U_A} = 0
  \qquad\Longrightarrow\qquad
  \pdv{S_A}{U_A} + \pdv{S_B}{U_A} = 0,
\]
and since $dU_A = -dU_B$ (total energy is conserved),
\[
  \pdv{S_A}{U_A} = \pdv{S_B}{U_B}.
\]
This is the equilibrium condition: both solids must have the same value of $\partial S/\partial U$.

\begin{center}
[DIAGRAM: Graph of $S_A$, $S_B$, and $S_{\text{total}}$ as functions of $q_A$. $S_A$ increases
monotonically, $S_B$ decreases monotonically, and $S_{\text{total}}$ has a smooth maximum. The
equilibrium point (marked with a cross) is where the slopes of $S_A$ and $S_B$ are equal in
magnitude. A steeper slope in $S_A$ than in $S_B$ at a given point means adding energy to $A$
gains more entropy than removing it from $B$ loses, so energy flows into $A$.]
\end{center}

The intuition is as follows. At a point where $\partial S_A/\partial U_A > \partial S_B/\partial
U_B$, transferring a small amount of energy from $B$ to $A$ increases the total entropy: the
entropy gained by $A$ exceeds the entropy lost by $B$. Energy will therefore flow from $B$ into
$A$. A steeper $\partial S/\partial U$ corresponds to a colder object (it gains more entropy per
unit of energy received), while a shallower slope corresponds to a hotter object (it gives away
energy more readily). This motivates the definition:

\begin{definition}
\textbf{Temperature} is defined by
\begin{equation}
  T \equiv \left(\pdv{S}{U}\right)_{N,V}^{-1} = \pdv{U}{S}\bigg|_{N,V}. \tag{1}
\end{equation}
Equivalently, $dS = dU/T$ at constant $N$ and $V$.
\end{definition}

% ---- Page 53: Real-World Examples ----

\subsection{Real-World Examples}

\subsubsection*{Einstein Solid (Harmonic Oscillators)}

Recall that for an Einstein solid in the high-temperature limit ($q \gg N$), the entropy is
\[
  S = Nk\ln\!\left(\frac{q}{N}\right) + Nk = Nk\ln U - Nk\ln(N\varepsilon) + Nk,
\]
where $U = q\varepsilon$. Differentiating with respect to $U$ at fixed $N$,
\[
  T = \left(\pdv{S}{U}\right)^{-1} = \frac{U}{Nk}
  \qquad\Longrightarrow\qquad U = NkT.
\]
This is the equipartition result for $N$ harmonic oscillators (each with two quadratic degrees
of freedom, giving $\frac{1}{2}kT$ per degree of freedom). No additional constants appear in
this formula---the definition of temperature in equation~(1) reproduces the correct physics
without any free parameters.

\subsubsection*{Monatomic Ideal Gas}

For a monatomic ideal gas the entropy from the Sackur--Tetrode equation depends on $U$ as
\[
  S = Nk\ln V + \frac{3}{2}Nk\ln U + \text{(function of }N\text{ only)}.
\]
Differentiating,
\[
  T = \left(\pdv{S}{U}\right)^{-1} = \frac{U}{\tfrac{3}{2}Nk}
  \qquad\Longrightarrow\qquad U = \frac{3}{2}NkT,
\]
which gives 3 degrees of freedom per atom (the three translational directions), each
contributing $\frac{1}{2}kT$ of energy. This is exactly the equipartition theorem for a
monatomic ideal gas. The factor $\frac{3}{2}$ reflects 3 degrees of freedom per oscillator.

% ---- Page 54: Problems (blank) ----

\subsection*{Problems}

% (No problems recorded on this page.)

% ---- Page 55: Section 3.2 — Entropy and Heat: Predicting Heat Capacities ----

\section{Entropy and Heat}

\subsection{Predicting Heat Capacities}

To compare statistical-mechanical predictions with experiment, we can use the
\textbf{heat capacity at constant volume}:
\begin{definition}
\[
  C_V \equiv \left(\pdv{U}{T}\right)_{N,V}.
\]
\end{definition}

For an Einstein solid in the high-temperature limit, $U = NkT$ (shown above), so $C_V =
\pdv{U}{T} = Nk$. For a monatomic ideal gas, $U = \frac{3}{2}NkT$, so
\[
  C_V = \pdv{T}\!\left(\frac{3}{2}NkT\right) = \frac{3}{2}Nk.
\]

\subsubsection*{Full Procedure to Find Heat Capacity}

The general statistical-mechanical program for computing $C_V$ proceeds as follows.
\begin{enumerate}[label=\arabic*.]
  \item Use quantum mechanics and combinatorics to find the multiplicity $\Omega$ in terms of
        $U$, $V$, $N$, and any other relevant variables.
  \item Take $S = k\ln\Omega$ to find the entropy.
  \item Differentiate $S$ with respect to $U$ to find $T(U, V, \ldots)$.
  \item Solve for $U$ as a function of $T$.
  \item Find $C_V \equiv \partial U/\partial T$.
\end{enumerate}

\subsection{Measuring Entropies}

For most systems it is too hard to write down an explicit formula for the multiplicity and
hence for the entropy. However, the entropy can be \emph{measured} thermodynamically. If a
small amount of heat $\dbar Q$ is added to a system while its volume is held constant (so no
work is done), then by the first law $dU = \dbar Q$, and the definition $T = (\partial S/\partial
U)^{-1}$ gives

\begin{formula}
\[
  dS = \frac{\dbar Q}{T} \qquad (\text{constant }V).
\]
\end{formula}

If the temperature of the object remains approximately constant while heat is added (or if
changes in temperature are tracked continuously), we can integrate:
\begin{formula}
\[
  \Delta S = S_f - S_i = \int_{T_i}^{T_f} \frac{C_V}{T}\,dT.
\]
\end{formula}
This formula allows the entropy of a system to be determined from calorimetric measurements of
$C_V(T)$ alone, without any microscopic model.

% ---- Page 56: Third Law, Macroscopic View of Entropy ----

\subsection{Third Law of Thermodynamics}

\begin{theorem}
\textbf{Third Law of Thermodynamics.} The entropy of a system approaches zero as its
temperature approaches absolute zero. The heat capacity $C_V$ also goes to zero as $T \to 0$.
\end{theorem}

In practice, $S$ is often non-zero at low temperatures because the system has \emph{residual
entropy}: degrees of freedom that remain ``frozen in'' and do not equilibrate on laboratory
time scales. Residual entropy can arise from frozen molecular arrangements, mixtures of nuclear
isotopes, multiplicity of alignment of nuclear spins, and similar effects. The Third Law
strictly applies only to equilibrium states.

Furthermore, as $T \to 0$, the heat capacity $C_V \to 0$ as well, because the various quantum
degrees of freedom (rotational, vibrational, etc.) ``freeze out'' one by one; they can no
longer be thermally excited and therefore stop contributing to the energy.

\subsection{Macroscopic View of Entropy}

The relation $dS = \dbar Q/T$ was historically the \emph{original} definition of entropy,
introduced by Clausius long before the statistical-mechanical interpretation. One can think of
energy as a fluid that is conserved as it flows between systems: energy is neither created nor
destroyed. Entropy, however, behaves differently: each ``packet'' of energy carries a little
entropy with it, and that entropy is \emph{not} conserved. It is \emph{created} whenever heat
flows (it is created wherever heat flows irreversibly). The total entropy of the universe
therefore never decreases.

% ---- Page 57: Section 3.3 — Paramagnetism ----

\section{Paramagnetism}

\subsection{Notation and Macroscopic Physics}

We now study a \textbf{two-state paramagnet}: a system of $N$ spin-$\tfrac{1}{2}$ dipoles
immersed in a constant external magnetic field $B$ in the $+\hat{z}$ direction. We assume
there are no interactions between dipoles --- this is the \emph{ideal} paramagnet.

Each dipole can point either ``up'' ($+\hat{z}$) or ``down'' ($-\hat{z}$), with energies
\[
  \varepsilon_{\uparrow} = -\mu B \quad (\text{``up'', lower energy, aligned with }B), \qquad
  \varepsilon_{\downarrow} = +\mu B \quad (\text{``down'', higher energy, anti-aligned with }B).
\]
Here $\mu$ is the magnitude of the magnetic moment of a single dipole. Let $N_\uparrow$ be the
number of up-spins and $N_\downarrow = N - N_\uparrow$ the number of down-spins. The total
energy is
\[
  U_{\text{total}} = \mu B(N_\uparrow - N_\downarrow) \cdot(-1) = -\mu B(N_\uparrow - N_\downarrow)
  = \mu B(N - 2N_\uparrow).
\]

\begin{definition}
The \textbf{magnetization} $M$ is the total magnetic moment of the entire system, with
``up'' contributing $+\mu$ and ``down'' contributing $-\mu$:
\[
  M \equiv \mu(N_\uparrow - N_\downarrow) = -\frac{U}{B}.
\]
\end{definition}

The multiplicity of the system depends only on $N_\uparrow$ (equivalently, on the energy or
the magnetization):
\begin{formula}
\[
  \Omega(N_\uparrow) = \binom{N}{N_\uparrow} = \frac{N!}{N_\uparrow!\,N_\downarrow!}.
\]
\end{formula}

% ---- Pages 58–65: Unit 1 Introduction — Statistical Description of Nature ----

\chapter{Introduction: Why Statistical Mechanics?}

\section{Motivation and Physical Examples}

The central question of statistical mechanics is: \emph{how does the macroscopic world arise
from the underlying microscopic one?} There is a vast gap between the level of description
at which we know the fundamental laws (Newton's equations for each particle, or Schr\"odinger's
equation for the quantum state) and the level at which we actually observe phenomena. Statistical
mechanics bridges that gap.

A partial list of systems where this question is non-trivial illustrates the breadth of the
problem:
\begin{itemize}
  \item A piston filled with gas
  \item Water in a kettle
  \item Air in the lungs
  \item A neutron star
  \item A fermionic superfluid
  \item A photon gas
\end{itemize}

\subsection{Efficiency of a Refrigerator}

A refrigerator does work to move heat from a cold reservoir (the food inside) to a hot
reservoir (the room). For an ideal refrigerator operating between room temperature $T_{\rm room}$
and food temperature $T_{\rm food}$, the ratio of heat removed to work done is bounded by:

\begin{center}
[DIAGRAM: Box labelled ``Fridge'' with an arrow labelled ``Do work'' entering and an arrow
labelled ``Heat out'' exiting to the right. This schematically represents the refrigeration
cycle: work is consumed to pump heat from cold to hot.]
\end{center}

\begin{formula}
\[
  \frac{\text{Heat removed}}{\text{Work done}} \leq \frac{T_{\rm room}}{T_{\rm room} - T_{\rm food}}
  \approx 18
\]
\end{formula}
for typical values. No matter how efficient your fridge is engineered to be, this ratio cannot
exceed the thermodynamic bound.

\subsection{Cooling a Quantum Gas}

\begin{center}
[DIAGRAM: Two columns showing $^3\mathrm{Li}$ (fermions) and $^6\mathrm{Li}$ (bosons). As the gas
is cooled (temperature decreasing downward), the energy level occupancy changes: bosons
condense into the ground state (Bose--Einstein condensation, marked BEC), while fermions stack
up in a Fermi sea due to the Pauli exclusion principle.]
\end{center}

Cooling a gas of quantum particles leads to qualitatively different behaviour depending on
whether the particles are bosons or fermions. Bosons can all occupy the same quantum state and
form a Bose--Einstein condensate (BEC), while fermions must fill up energy levels one by one
from the bottom.

\subsection{Superfluidity}

\begin{center}
[DIAGRAM: A flat plane (representing a superfluid film) confined in a geometry. On the left,
cooling causes the system to become a superfluid; on the right, the superfluid exhibits
quantum vortices---topological defects in the phase of the order parameter.]
\end{center}

A superfluid is a phase of matter in which a quantum fluid flows without viscosity. A remarkable
feature of rotating superfluids is the appearance of \emph{quantum vortices}: each vortex
carries exactly one quantum of circulation, $h/m$. This is a macroscopic manifestation of
quantum mechanics.

\section{Why Do We Need a Statistical Description?}

\subsection{The Complexity Argument}

The fundamental obstacle to a direct microscopic approach is one of scale. Consider air in an
ordinary room:
\begin{itemize}
  \item $V \sim 10^3\,\mathrm{m}^3$
  \item $P \approx 1\,\mathrm{atm} \approx 10^5\,\mathrm{Pa}$
  \item $T \approx 300\,\mathrm{K}$
\end{itemize}
From the ideal gas law $PV = Nk_BT$,
\[
  N \approx 10^{28} \text{ molecules}.
\]
To specify the full state of the gas one would need the position and velocity of every molecule:
$(x, y, z, v_x, v_y, v_z)$ for each, giving $6N \sim 6 \times 10^{28}$ numbers. Storing these
as single-precision floats at $10^{16}$ bytes per hard drive would require $10^{20}$~TB of
storage. Even on the fastest current supercomputer (performing $\sim 10^{18}$ floating-point
operations per second), a single update of the state would take on the order of $10^{10}$
seconds --- far longer than the age of the universe.

But the deeper obstacle is not merely computational. Even if we could solve Laplace's demon's
problem and track every particle, many of the most important macroscopic phenomena simply do
not \emph{appear} in the fundamental microscopic laws:
\begin{itemize}
  \item The concepts of \emph{temperature} and \emph{entropy} are nowhere in Newton's equations.
  \item $F = ma$ says nothing about the arrow of time, phase transitions, or thermal
        equilibration.
\end{itemize}
As Philip Anderson put it: \textbf{More is Different}. As the complexity of a system increases,
qualitatively new concepts and laws emerge that cannot be reduced to the laws governing the
constituents:
\[
  \text{Quarks/electrons} \;\to\; \text{Atoms/molecules} \;\to\; \text{Materials} \;\to\;
  \text{Engines.}
\]
The basic idea of statistical mechanics is to turn this complexity to our advantage by adopting
a statistical model. Even if the exact microscopic laws are unknown, we can still make powerful
predictions by characterising the probability distribution over microstates.

\subsection{Statistical Mechanics vs.\ Thermodynamics}

The two pillars of the subject complement each other:

\begin{center}
\begin{tabular}{p{0.45\linewidth}|p{0.45\linewidth}}
  \textbf{Statistical Mechanics} & \textbf{Thermodynamics} \\[4pt]
  Connects microscopic models to macroscopic observables. & Governs the flow of energy, heat,
  and work at the macroscopic level. \\[4pt]
  Inputs: atoms, spins, EM modes, \ldots & Language of emergent properties. \\[4pt]
  Calculates statistical behaviour. & Temperature $\sim$ average thermal energy; pressure $\sim$
  average force on wall. \\[4pt]
  & Thermodynamics does not concern itself with the microscopic model. \\[4pt]
  & Abstractness is the price of generality.
\end{tabular}
\end{center}

\section{Some Tools We Will Need}

The main mathematical and conceptual tools used throughout this course are:
\begin{itemize}
  \item \textbf{Quantum mechanics} --- statistical mechanics requires knowing the states of the
        system. Quantum mechanics provides a discrete, countable set of states (energy levels
        labelled by quantum numbers), which makes probability theory well-defined.
  \item \textbf{Probability theory} --- used to calculate averages, fluctuations, and the
        behaviour of the system in equilibrium.
  \item \textbf{Partial derivatives} --- used to relate changes in thermodynamic quantities.
  \item \textbf{Some base concepts} --- introduced below.
\end{itemize}

\section{Basic Concepts}

\subsection{Thermodynamic System and Environment}

A \textbf{thermodynamic system} consists of many constituents (molecules, spins, photons,
etc.). The \textbf{environment} (or \textbf{reservoir}) is the rest of the universe. The system
may exchange energy, volume, or particles with its environment, depending on the constraints
imposed.

\subsection{Internal Energy}

\begin{definition}
The \textbf{internal energy} $U$ of a system is the sum of all contributions to the energy of
the system treated as an isolated whole:
\[
  U := \sum_i E_i.
\]
\end{definition}

For example, for a box of helium, $U$ includes the kinetic energy of the atoms, the
electron--nucleus interactions (binding energies), and the nuclear binding energies. For a box
of oxygen, it additionally includes rotational and vibrational energies of the $\mathrm{O}_2$
molecules.

Internal energy does \emph{not} include the centre-of-mass kinetic energy of the entire box
(bulk motion), nor the potential energy of the system in an external field (e.g.\ gravity):
those are not intrinsic to the internal state of the system.

\subsection{Temperature}

\begin{center}
[DIAGRAM: Graph of internal energy $U$ vs.\ temperature $T$. The curve is flat at $T = 0$ and
rises steeply at higher $T$. At $T = 0$, no more energy can be extracted unless a chemical
reaction occurs (marked ``chemical reaction''). The horizontal axis shows $T = 0$, $T_{\rm ice}$,
and $T_{\rm water}$ as reference points.]
\end{center}

Temperature is a surprisingly elusive concept. The boxed statement below captures the
operational meaning:
\begin{fact}
$T = 0$ means that no more energy can be extracted from the system (kinetic energy $= 0$)
unless a chemical reaction occurs.
\end{fact}

We use the \textbf{Kelvin scale}: $0^\circ\mathrm{C} = 273.15\,\mathrm{K}$ and a change of
$1^\circ\mathrm{C}$ equals a change of $1\,\mathrm{K}$.

\subsection{Thermal Energy}

\begin{definition}
The \textbf{thermal energy} is $U(T) - U(0)$, i.e.\ the internal energy measured relative to
the ground state. It differs from $U$ only by a (sometimes inconvenient) constant.
\end{definition}

\subsection{Heat}

\begin{definition}
\textbf{Heat} $Q$ is the piece of thermal energy transferred between systems at the microscopic
(molecular) level. Heat flows from hot to cold. Heat is \emph{not} a type of energy stored in
a system; it is a \emph{mode of energy transfer}, not a property of a system.
\end{definition}

\subsection{Thermal Equilibrium}

\begin{definition}
\textbf{Thermal equilibrium}: if an isolated system is left alone, its macroscopic properties
cease to change over time. Temperatures equalize and there is no further net heat flow.
\end{definition}

\section{Zeroth Law of Thermodynamics}

\begin{theorem}
\textbf{Zeroth Law of Thermodynamics.} If system $A$ is in thermal equilibrium with system $B$,
and system $B$ is in thermal equilibrium with system $C$, then $A$ is in thermal equilibrium
with $C$.
\end{theorem}

This law is what makes temperature a well-defined, transitive property: one can always use a
thermometer (system $B$) to compare the temperatures of two other systems.

\section{State of a System}

The full description of a system differs markedly between classical and quantum mechanics.

\begin{center}
\begin{tabular}{p{0.45\linewidth}|p{0.45\linewidth}}
  \textbf{Classical} & \textbf{Quantum} \\[4pt]
  Each particle $i$ has a position $\vb{r}_i$ and velocity $\vb{v}_i$. The state is the
  $6N$-dimensional vector $\{\vb{r}_1, \vb{v}_1, \vb{r}_2, \vb{v}_2, \ldots, \vb{r}_N, \vb{v}_N\}$.
  & Only discrete energies are allowed. States are labelled by quantum numbers, e.g.\
  $n = 1, 2, 3, \ldots$. A quantum state $\ket{\psi}$ is a list of quantum numbers specifying
  the occupation of each level.
\end{tabular}
\end{center}

\begin{center}
[DIAGRAM: Classical side: $N$ particles with arrows indicating positions and velocities.
Quantum side: a ladder of energy levels $n = 1, 2, \ldots, k, \ldots$ with some levels marked
as occupied (countable).]
\end{center}

\subsection{Quantum States (Continued)}

A concrete example of a quantum state is a collection of elementary magnets (spins): each may
be $\uparrow$ ($+1$) or $\downarrow$ ($-1$). A possible state of the system is
\[
  \ket{\psi} = \{\tfrac{1}{2}, +1, +1, -1, +1, -1, \ldots\}.
\]
This is a list of quantum numbers, one per spin. Both classical and quantum descriptions aim to
fully specify the \textbf{microstate} of the system. In practice, even specifying the microstate
of $10^{28}$ molecules is either impossible or useless --- so we must work statistically.

\section{States in Thermodynamics}

\subsection{Reaching Equilibrium}

In practice, we allow the system to ``settle,'' meaning it evolves until it ceases to change
at the macroscopic level (e.g.\ temperature $T$, pressure $P$, volume $V$, and energy $U$ all
become time-independent). The system is then said to have reached \textbf{thermodynamic
equilibrium}. It may not be in any specific microstate; the microstate fluctuates, but the
macrostate is fixed.

\subsection{Macrostates}

\begin{definition}
A \textbf{macrostate} is defined by \emph{state variables} --- macroscopic quantities that
characterise the equilibrium state. Examples include $N$ (number of particles), $V$ (volume),
$U$ (internal energy), $P$ (pressure), $T$ (temperature), and $S$ (entropy). These are not
all independent.
\end{definition}

One can fully characterise a macrostate by a set of \emph{independent} state variables. For
example, specifying $(U, N, V)$ for an ideal gas determines all other equilibrium properties.
Some state variables are defined even away from equilibrium (e.g.\ $N$), while others are not.

\subsubsection*{Extensive and Intensive Variables}

State variables are classified by how they scale with system size.
\begin{itemize}
  \item \textbf{Extensive} variables scale with $N$ (system size): $N$, $V$, $U$, $S$. If two
        identical systems are combined, an extensive variable doubles.
  \item \textbf{Intensive} variables are independent of system size: $T$, $P$. If two identical
        systems are combined, an intensive variable is unchanged.
\end{itemize}

For example, combining two identical systems $\boxed{1}$ and $\boxed{2}$:
\[
  N = N_1 + N_2 = 2N \;\to\; \text{Extensive}, \qquad
  V = V_1 + V_2 = 2V \;\to\; \text{Extensive},
\]
\[
  T = T_1 = T_2 \;\to\; \text{Intensive}, \qquad
  P = P_1 = P_2 \;\to\; \text{Intensive}, \qquad
  S = S_1 + S_2 \;\to\; \text{Extensive}.
\]

\subsection{State Functions}

\begin{definition}
A \textbf{state function} is a quantity whose value depends only on the current macrostate, not
on the history of how that state was reached. Formally, $f$ is a state function when
\[
  \pdv{f}{y}\bigg|_z = \pdv{f}{z}\bigg|_y \quad (\text{the mixed partials commute}).
\]
\end{definition}

The macrostate of a thermodynamic system may itself depend on the choice of state variables:
for example, the phase diagram of a substance (solid, liquid, gas) changes depending on which
variables $(P, T)$ are used. Phase diagrams and $P$-$T$ diagrams (with regions for solid,
liquid, gas, and exotic phases like superconductors, antiferromagnets, and strange metals)
illustrate how richly the macrostate can vary.

\begin{center}
[DIAGRAM: Left panel: ordinary $P$-$T$ phase diagram showing solid, liquid, and gas regions
with the triple point and critical point. Right panel: a more exotic phase diagram (e.g.\ for
a high-temperature superconductor) showing regions labelled AFM (antiferromagnet), strange
metal, $\mathrm{S.C.}$ (superconductor), and Fermi liquid, with ``doping'' on the horizontal
axis and temperature on the vertical axis.]
\end{center}

\section{Microstates and the Information Gap}

\begin{definition}
A \textbf{microstate} defines the specific microscopic configuration of all the constituents
of a system.
\begin{itemize}
  \item \textbf{Classical}: specified by $6N$ numbers $\{\vb{r}_i, \vb{v}_i\}$.
  \item \textbf{Quantum}: specified by $\sim e^N$ quantum numbers $\{n_1, n_2, \ldots\}$ (one
        per accessible level).
\end{itemize}
\end{definition}

One can uniquely identify a microstate with a point in \textbf{phase space} (the
$6N$-dimensional space of positions and momenta). For a single one-dimensional particle, the
phase space is two-dimensional (position $x$ and momentum $p$). For $N$ particles in $d$
dimensions, the phase space has dimension $D = 2Nd$.

\begin{fact}
\textbf{Crucial Information Gap.} There is a vast gap between the macrostate and the
microstate. Every macrostate is compatible with an enormous number of microstates --- the
\emph{multiplicity} $\Omega$. Macroscopic properties depend on the macrostate (e.g.\ heat
capacity, compressibility), not on which particular microstate the system happens to be in.
\end{fact}

\textbf{Example.} A box of gas with macrostate $(N, V, U)$ is compatible with $\Omega(N, V, U)$
microstates. The multiplicity is the key quantity linking the microscopic and macroscopic
descriptions.

\begin{definition}
The \textbf{micro-canonical ensemble} is the set of all copies (``replicas'') of the
microstate of a system with fixed macrostate $(N, V, U)$. It is the natural ensemble for an
isolated system.
\end{definition}
